\section{Future Directions}
\label{sec:future_directions}

Table~\ref{tab:survey_comparison} clarifies the contribution of this
survey relative to prior work by comparing its scope with five recent
reviews of Speech LLMs and audio-language models. 

The remainder of this
section adopts a lifecycle perspective to organize future directions
around stage-specific challenges and cross-stage dependencies that shape
practical Speech LLM systems.

\begin{table*}[t]
\caption{Coverage comparison: this survey versus five recent Speech LLM
and audio-language reviews. \checkmark{} indicates substantive coverage;
\textendash{} indicates limited or absent coverage. In the last row,
qualifiers in italics identify dimensions where this survey adds
coverage not present in the earlier reviews.}
\label{tab:survey_comparison}
\centering
\footnotesize
\begin{tabularx}{\textwidth}{p{0.18\textwidth}XXXXXX}
\toprule
\textbf{Survey} & \textbf{Training} & \textbf{Meta-evaluation of benchmarks} &
\textbf{Deployment / serving} & \textbf{Safety / privacy} &
\textbf{Lifecycle coupling} & \textbf{Methodology type} \\
\midrule
Cui et al.\ 2025 \citep{cui2024recent} &
\checkmark{} & \textendash{} & \textendash{} & \textendash{} &
\textendash{} & Narrative \\
\midrule
Peng et al.\ 2025 \citep{peng2025surveyspeechlargelanguage} &
\checkmark{} & Capability taxonomy only & \textendash{} &
\textendash{} & \textendash{} & Narrative \\
\midrule
Ji et al.\ (WavChat) \citep{ji2024wavchat} &
Partial & \textendash{} & Dialogue / duplex only & \textendash{} &
\textendash{} & Narrative \\
\midrule
Latif et al.\ 2023 \citep{latif2023sparks} &
\checkmark{} (speech foundation models) & \textendash{} &
\textendash{} & \textendash{} & \textendash{} & Narrative \\
\midrule
Ye et al.\ 2024 (LVLM safety) \citep{ye2024survey} &
\textendash{} & \textendash{} & \textendash{} &
\checkmark{} (vision-centric) & \textendash{} & Narrative \\
\midrule
\textbf{This survey} &
\checkmark{} & \checkmark{} \emph{(validity-aware)} &
\checkmark{} \emph{(interaction-oriented)} &
\checkmark{} \emph{(speech-specific)} &
\checkmark{} \emph{(cross-stage)} & Lifecycle audit \\
\bottomrule
\end{tabularx}
\end{table*}

\paragraph{Training} Closing the data imbalance for low-resource
languages and non-Latin scripts will require automated data generation,
cross-lingual transfer, and quality-aware weakly supervised pipelines.
Alignment modules need to preserve paralinguistic content while
remaining streaming-friendly, and preference alignment needs reward
signals that capture acoustic and prosodic quality, not only transcript
fidelity.

\paragraph{Evaluation} Benchmarks should adopt joint evaluation across
content fidelity, paralinguistic awareness, interaction dynamics,
latency, and safety; distribution-aware metrics for subjective targets;
mandatory human baselines for new capability claims; and standardized
Speech-LLM Evaluation Cards for cross-paper comparability
(Section~\ref{sec:evaluation_benchmarking}).

\paragraph{Deployment} Compositional optimization---combining
quantization, speculative decoding, KV-cache reduction, and streaming
serving in a single full-duplex stack---remains an open engineering
problem. Reporting standards should include latency under simultaneous
input and output, not only offline batch throughput
(Section~\ref{sec:deployment}).

\paragraph{Safety, privacy, and governance} Audio-specific red teaming,
voice anonymization that resists adaptive purification, watermarking
that survives realistic transformations, and modality-aware governance
documentation are all immediate priorities
(Section~\ref{sec:safety_privacy_governance}).

\paragraph{Cross-stage coupling} The most impactful research questions
sit between stages: training pipelines that anticipate streaming
deployment, benchmarks that score safety and capability jointly, and
governance frameworks that translate technical risk into deployable
controls. A lifecycle perspective, rather than further specialization
within a single stage, is likely to produce the largest gains for
real-world Speech LLM systems.
